\documentclass[11pt]{article}

\usepackage[margin=0.85in]{geometry}
\usepackage{booktabs}
\usepackage{array}
\usepackage{tabularx}
\usepackage{multirow}
\usepackage{xcolor}
\usepackage{colortbl}
\usepackage{siunitx}
\usepackage{enumitem}
\usepackage{hyperref}

\hypersetup{
  colorlinks=true,
  linkcolor=blue!50!black,
  urlcolor=blue!50!black,
  citecolor=blue!50!black
}

\definecolor{headerblue}{HTML}{1F4E79}
\definecolor{softblue}{HTML}{EAF2F8}
\definecolor{softgreen}{HTML}{EAF7EA}
\definecolor{softgold}{HTML}{FFF4D6}
\definecolor{softgray}{HTML}{F4F6F7}
\definecolor{darkgray}{HTML}{3A3A3A}

\newcommand{\best}[1]{\cellcolor{softgreen}\textbf{#1}}
\newcommand{\warn}[1]{\cellcolor{softgold}#1}
\newcommand{\metric}[1]{\textbf{#1}}

\sisetup{
  round-mode=places,
  round-precision=4,
  detect-weight=true,
  detect-family=true
}

\title{\vspace{-1.0em}\textbf{InvariantForecasting: First Run Documentation}\\[-0.15em]
\large Base vs. GRPO $\lambda=0$ vs. GRPO $\lambda=1$ on ForecastBench Current}
\author{Internal experiment log}
\date{April 28, 2026}

\begin{document}
\maketitle
\vspace{-1.5em}

\section*{Executive Summary}

This document records the first end-to-end experimental run for the
\texttt{InvariantForecasting} project. The goal was to test whether GRPO
post-training with an explicit frame-invariance penalty reduces paraphrase
variance in LLM forecasts relative to both the base model and outcome-only GRPO.

\textbf{Post-run caveat.} A later audit found that this run used prompts where
many ForecastBench template dates were still literal placeholders such as
\texttt{\{resolution\_date\}} and \texttt{\{forecast\_due\_date\}}, and the
paraphrase artifact did not retain full background and resolution context. The
results below should therefore be treated as a debugging baseline rather than a
clean research result.

The main empirical finding is that GRPO improves parseable forecast coverage and
reduces paraphrase variance. The explicit invariance run ($\lambda=1$) gives the
lowest paraphrase standard deviation while slightly improving Brier relative to
the outcome-only GRPO run ($\lambda=0$). The base model still has the best Brier
among parseable groups, but its coverage is much lower, so that accuracy number
is conditional on a smaller subset of outputs.

\section*{Result Table}

\begin{table}[h!]
\centering
\scriptsize
\renewcommand{\arraystretch}{1.28}
\setlength{\tabcolsep}{3.5pt}
\rowcolors{3}{softgray}{white}
\begin{tabular}{llrrrrrr}
\rowcolor{headerblue}
\color{white}\textbf{Run} &
\color{white}\textbf{Objective} &
\color{white}\textbf{Brier $\downarrow$} &
\color{white}\textbf{ParaSD $\downarrow$} &
\color{white}\textbf{Var. Cov. $\uparrow$} &
\color{white}\textbf{Grp. Cov. $\uparrow$} &
\color{white}\textbf{N} &
\color{white}\textbf{Preds.} \\
\midrule
Base Qwen-7B & No RL & \best{0.2436} & 0.1229 & 0.6667 & 0.2833 & 60 & 300 \\
GRPO $\lambda=0$ & Brier only & 0.2508 & 0.0894 & \best{1.0000} & \best{1.0000} & 60 & 300 \\
GRPO $\lambda=1$ & Brier + variance & 0.2495 & \best{0.0800} & \best{1.0000} & \best{1.0000} & 60 & 300 \\
\bottomrule
\end{tabular}
\caption{Held-out test results on 60 ForecastBench-current question groups, with 5 semantically equivalent variants per group. Lower Brier is better. Lower ParaSD means less variation across paraphrases. Coverage is the fraction of outputs from which a strict probability forecast could be parsed.}
\label{tab:first-run-results}
\end{table}

\section*{Key Comparisons}

\begin{itemize}[leftmargin=*]
  \item \textbf{Parseability improved substantially.} Variant-level coverage increased from 0.6667 for the base model to 1.0000 for both GRPO runs.
  \item \textbf{Paraphrase variance decreased.} Mean ParaSD dropped from 0.1229 for the base model to 0.0894 for $\lambda=0$ and 0.0800 for $\lambda=1$.
  \item \textbf{$\lambda=1$ improved invariance beyond outcome-only GRPO.} Relative to $\lambda=0$, $\lambda=1$ reduced ParaSD by 10.49\%.
  \item \textbf{$\lambda=1$ did not add an accuracy cost relative to $\lambda=0$.} Mean Brier improved slightly from 0.2508 to 0.2495.
  \item \textbf{Base Brier remains best, with an important caveat.} The base model's Brier of 0.2436 is computed on parseable groups only, while the GRPO runs reached complete variant coverage. This means the base result is not a full-coverage forecast result.
\end{itemize}

\section*{Numerical Deltas}

\begin{table}[h!]
\centering
\renewcommand{\arraystretch}{1.2}
\rowcolors{3}{softgray}{white}
\begin{tabular}{lrrrr}
\rowcolor{headerblue}
\color{white}\textbf{Comparison} &
\color{white}\textbf{$\Delta$ Brier} &
\color{white}\textbf{$\Delta$ Brier \%} &
\color{white}\textbf{ParaSD Reduction} &
\color{white}\textbf{ParaSD Reduction \%} \\
\midrule
$\lambda=0$ vs. base & +0.0072 & +2.96\% & 0.0335 & 27.24\% \\
$\lambda=1$ vs. base & +0.0058 & +2.39\% & 0.0428 & 34.87\% \\
$\lambda=1$ vs. $\lambda=0$ & -0.0014 & -0.55\% & 0.0094 & 10.49\% \\
\bottomrule
\end{tabular}
\caption{Absolute and relative changes from the first run. Positive Brier deltas are worse; ParaSD reductions are better.}
\end{table}

\section*{Experimental Setup}

\subsection*{Dataset}

The run used the \texttt{forecastbench/current} data assembled for this project,
filtered to resolved binary markets with resolution dates after August 31, 2025.
Each original question was paired with 5 semantically equivalent variants: the
original wording plus four generated paraphrases. The resulting paraphrased file
contains 598 question groups and 2,990 variant rows.

The code used a deterministic group-level split:

\begin{center}
\begin{tabular}{lrr}
\toprule
\textbf{Split} & \textbf{Groups} & \textbf{Variant Rows} \\
\midrule
Train & 478 & 2,390 \\
Validation & 60 & 300 \\
Test & 60 & 300 \\
\bottomrule
\end{tabular}
\end{center}

Only the held-out test split was used for the headline results in Table
\ref{tab:first-run-results}. We did not report all 598 groups as the main result
because that would mix training, validation, and test questions after GRPO
post-training.

\subsection*{Model}

The base model was \texttt{Qwen/Qwen2.5-7B-Instruct}. The GRPO runs used LoRA
post-training with rank 32, alpha 64, dropout 0.05, and attention projection
modules as targets: \texttt{q\_proj}, \texttt{k\_proj}, \texttt{v\_proj}, and
\texttt{o\_proj}.

\subsection*{Prompt and Parsing}

Each model received a forecasting prompt asking for a short rationale and a final
line formatted exactly as:

\begin{verbatim}
Probability: <number between 0 and 1>
\end{verbatim}

During debugging we found that a loose parser could incorrectly extract numbers
from rationales, list markers, dates, or prompt fragments. The final reported
results therefore use a strict parser that only accepts explicit
\texttt{Probability:} or \texttt{Final answer:} forecasts and uses the last valid
forecast if multiple appear.

\subsection*{Reward}

For a question group with paraphrase variants $x_1, \ldots, x_K$, model forecasts
$p_i$, binary outcome $Y \in \{0,1\}$, and consensus forecast
$\bar{p}=K^{-1}\sum_i p_i$, the reward was:

\[
R_i = -(p_i - Y)^2 - \lambda (p_i - \bar{p})^2.
\]

The $\lambda=0$ run is the outcome-only Brier reward baseline. The $\lambda=1$
run adds an explicit paraphrase-variance penalty.

\subsection*{Evaluation Metrics}

\begin{itemize}[leftmargin=*]
  \item \textbf{Mean Brier:} computed at the question-group level using the mean parsed probability across variants as the consensus forecast.
  \item \textbf{Mean ParaSD:} population standard deviation of parsed probabilities across paraphrases, averaged over question groups.
  \item \textbf{Variant coverage:} fraction of all 300 test variant outputs with a parseable probability.
  \item \textbf{Full-group coverage:} fraction of the 60 test groups where all 5 variants yielded parseable probabilities.
\end{itemize}

\section*{Step-by-Step Process Completed}

\begin{enumerate}[leftmargin=*]
  \item \textbf{Prepared the repository.} The \texttt{InvariantForecasting} repository was set up with data, scripts, configs, and training/evaluation modules.
  \item \textbf{Prepared ForecastBench-current data.} We used 598 resolved binary question groups after the August 31, 2025 cutoff, with 5 variants per group.
  \item \textbf{Validated grouped splits.} The loader confirmed 478 training groups, 60 validation groups, and 60 held-out test groups.
  \item \textbf{Hardened probability parsing.} The parser was changed to avoid extracting irrelevant numbers from model rationales.
  \item \textbf{Fixed decoder-only evaluation behavior.} Evaluation was patched to use left padding and to slice generated completions correctly.
  \item \textbf{Patched TRL compatibility.} The trainer now drops unsupported \texttt{GRPOConfig} arguments for the installed TRL version and sets evaluation batch sizing safely.
  \item \textbf{Protected grouped rewards.} The trainer now checks that paraphrase groups are not split across reward batches and that dataset shuffling is disabled.
  \item \textbf{Disabled intermediate validation for real runs.} Expensive mid-training validation was disabled with \texttt{eval\_strategy: "no"}; final held-out evaluation is still run after training.
  \item \textbf{Ran the strict base evaluation.} Base Qwen2.5-7B was evaluated zero-shot on the held-out test split using the strict parser.
  \item \textbf{Ran GRPO $\lambda=0$.} The model was post-trained with the Brier reward only, then evaluated on the held-out test split.
  \item \textbf{Ran GRPO $\lambda=1$.} The model was post-trained with Brier plus explicit paraphrase-invariance penalty, then evaluated on the held-out test split.
  \item \textbf{Exported result artifacts.} The relevant summaries, group metrics, and variant predictions were packaged for local analysis.
\end{enumerate}

\section*{Interpretation}

This first run supports the narrow claim that explicit invariance regularization
can reduce paraphrase sensitivity beyond outcome-only GRPO. The $\lambda=1$ run
achieved the best ParaSD, reducing paraphrase variance by 34.87\% relative to the
base model and by 10.49\% relative to $\lambda=0$. It also slightly improved
Brier relative to $\lambda=0$, suggesting that the invariance penalty did not add
an extra accuracy cost in this run.

The result does not yet support a strong claim that GRPO improves forecasting
accuracy over the base model, since the base model has a lower Brier among
parseable groups. Instead, the current finding is best framed as:

\begin{quote}
On this held-out ForecastBench-current split, outcome-only GRPO improved
parseability and reduced paraphrase variance, while adding an explicit
paraphrase-invariance penalty reduced variance further with no additional Brier
cost relative to outcome-only GRPO.
\end{quote}

\section*{Limitations and Next Steps}

\begin{itemize}[leftmargin=*]
  \item \textbf{Small held-out test set.} The current headline result uses 60 question groups. Bootstrap confidence intervals are needed.
  \item \textbf{Single seed.} The result should be repeated across at least 2--3 seeds.
  \item \textbf{Limited $\lambda$ sweep.} A useful next ablation is $\lambda \in \{0, 0.5, 1, 2, 5\}$.
  \item \textbf{Coverage asymmetry.} Base Brier is measured under lower parse coverage, while GRPO reaches full coverage. A parse-failure-aware metric should also be reported.
  \item \textbf{Possible 0.5 collapse.} We should inspect prediction distributions to confirm that lower ParaSD is not merely constant-output behavior.
  \item \textbf{Held-out paraphraser needed.} The next robustness test should evaluate on paraphrases generated by a different paraphrase method.
\end{itemize}

\section*{Commands Used for Final Evaluations}

\subsection*{Base Strict Evaluation}

\begin{verbatim}
PYTHONPATH=src python -m frame_invariance.training.evaluate \
  --config configs/training/grpo_forecastbench.yaml \
  --model Qwen/Qwen2.5-7B-Instruct \
  --run-name base_qwen7b_strict \
  --split test \
  --batch-size 4
\end{verbatim}

\subsection*{$\lambda=0$ Evaluation}

\begin{verbatim}
PYTHONPATH=src python -m frame_invariance.training.evaluate \
  --config configs/training/grpo_forecastbench.yaml \
  --model outputs/grpo_forecastbench_qwen7b_lambda0 \
  --base-model Qwen/Qwen2.5-7B-Instruct \
  --run-name grpo_lambda0 \
  --split test \
  --batch-size 4
\end{verbatim}

\subsection*{$\lambda=1$ Evaluation}

\begin{verbatim}
PYTHONPATH=src python -m frame_invariance.training.evaluate \
  --config configs/training/grpo_forecastbench.yaml \
  --model outputs/grpo_forecastbench_qwen7b_lambda1 \
  --base-model Qwen/Qwen2.5-7B-Instruct \
  --run-name grpo_lambda1 \
  --split test \
  --batch-size 4
\end{verbatim}

\section*{Artifact Paths}

The main result artifacts produced on RunPod were:

\begin{verbatim}
results/grpo_lambda0/test/summary.json
results/grpo_lambda0/test/group_metrics.csv
results/grpo_lambda0/test/variant_predictions.csv

results/grpo_lambda1/test/summary.json
results/grpo_lambda1/test/group_metrics.csv
results/grpo_lambda1/test/variant_predictions.csv
\end{verbatim}

The export archive was created as:

\begin{verbatim}
invariantforecasting_results.tar.gz
\end{verbatim}

\end{document}
